\section{Conclusion}
This paper introduces TinyNet, a novel routing protocol tailored for Networked Control Systems, emphasizing Message Delivery Time Determinism over Total Throughput. TinyNet's efficacy in routing messages using minimalistic hardware has been showcased alongside a simulation tool that assesses its performance. A standout feature of TinyNet is its ability to manage message routing in extensively connected networks without relying on a Network Controller or requiring nodes to possess complete network topology knowledge.

The pivotal contribution of this work is the creation of a port-forwarding strategy that utilizes real-time data about queue sizes at subsequent hops. This innovation allows for the intelligent rerouting of packets to circumvent congested network segments.

Determining realistic usage profiles for Network Control Systems has been challenging due to limited existing research. We plan to gather primary data on practical applications and challenges in this field by consulting industry stakeholders at companies like Boeing, Airbus, Kittyhawk, and Moog Inc.

The use of UART presents limitations in our project, specifically in capping the permissible bitrate to only 3.125 Mbps and necessitating a ``stateful'' environment where all ports must share the same baudrate settings. To address these issues, one future direction could include an FPGA-based link layer, implementing a ``co-clocking'' technique for establishing bitrate. In this approach, neither side of a transmission pre-defines a bitrate. Instead, data transfer occurs in a responsive manner: once one side shifts data into a register, it sends an acknowledgment bit, triggering the transmission of the next bit. This method has successfully achieved bitrates of up to 65 Mbps. We have integrated this system with a microcontroller analogous to the ATSAMS70 used in the TinyNet router, achieving similar transmission speeds via a 5-MHz wide parallel bus. This advancement aims to establish a completely stateless, configuration-free system, significantly enhancing the flexibility and scalability of TinyNet's network infrastructure.

In addition to leveraging FPGAs for a stateless link layer, FPGAs can also be used for switching and port-forwarding, leveraging the simplicity of the TinyNet protocol for implementation in Verilog. This effort is expected to surpass the performance of Switched Ethernet while preserving the stateless, fault-tolerant, and adaptive qualities of TinyNet.

Finally, the current routing protocol employs a lambda function, as outlined in equation (\ref{eq:cost}), for navigating around busy ports. We aim to explore machine learning techniques to refine this function, optimizing the routing process based on network topology and cost considerations, thereby enhancing overall network performance.